% TEMPLATE-MILC-v3.tex - plantilla maestra UNICA de las guias MILC v3.
%
% La usan las tres familias de guias, cada una con su driver:
%   · Grados 6-11  -> scripts/build-guias-g11.py
%   · Bebras       -> scripts/build-guias-bebras.py
%   · Semillero    -> scripts/build-guias-semillero.py
%
% Antes habia tres copias casi identicas (~400 lineas iguales, 6 distintas):
% cada mejora habia que aplicarla tres veces, y en la practica no se hacia
% (el motor de recursos vivia solo en dos de ellas). Lo que varia entre
% familias esta parametrizado en 6 placeholders que cada driver llena
% (se nombran aqui SIN su sintaxis real: el builder reemplaza por texto plano,
% tambien dentro de los comentarios, y un valor multilinea rompe el preambulo):
%
%   HEADER_IZQ        encabezado izquierdo de pagina
%   HEADER_DER        encabezado derecho de pagina
%   PORTADA_ETIQUETA  rotulo sobre el numero grande (GRADO/MOMENTO/MODULO)
%   PORTADA_SERIE     linea de serie de la portada
%   PORTADA_ID        identificador de la guia en la portada
%   PORTADA_CIRCULO   el circulito de grado (°), o vacio si no aplica
%
% El resto de claves las documenta content/guias/_SCHEMA.md. El script valida
% que no queden placeholders sin reemplazar antes de escribir el archivo final.

\documentclass[11pt,letterpaper]{article}
\usepackage[margin=2.0cm]{geometry}
\usepackage[table]{xcolor}
\usepackage{fontspec}
\usepackage[spanish]{babel}
\usepackage{tikz}
% Librerías TikZ para los diagramas de las guías (bloque `recursos:`):
%  · babel  → babel en español vuelve activos caracteres como «>», y TikZ los usa
%             en sus opciones (>=stealth, ->). Sin esta librería, cualquier
%             diagrama con flechas revienta con «\language@active@arg> has an
%             extra }». Debe ir ANTES de que se dibuje nada.
%  · positioning        → sintaxis «below left=2pt and 3pt».
%  · decorations.pathreplacing → llaves ({brace}) para señalar rangos.
\usetikzlibrary{babel,positioning,decorations.pathreplacing}
\usepackage{graphicx}
% Circuitos: el trabajo de electrónica se hace hoy con micro:bit y sus sensores
% integrados (MakeCode, sin cableado), así que no se necesitan esquemáticos.
% Si algún día entran montajes con protoboard, basta descomentar esta línea:
% los diagramas .tex/.tikz declarados en `recursos:` ya se incrustan solos.
% \usepackage{circuitikz}
\usepackage[most]{tcolorbox}
\usepackage{tabularx}
\usepackage{array}
\usepackage{fancyhdr}
\usepackage{titlesec}
\usepackage[document]{ragged2e}  % bandera a la derecha en todo el cuerpo
\usepackage{enumitem}
\usepackage{microtype}

% Helvetica Neue no trae la flecha U+2192, asi que cada «→» escrito literalmente
% en el YAML se compilaba a NADA, en silencio (462 flechas en 61 guias: rutas del
% tipo «paleta → tipografia → lockup» salian rotas). Mapeamos el caracter a su
% version matematica, que si tiene glifo. Asi el autor puede seguir escribiendo
% «→» comodamente en el YAML.
\usepackage{newunicodechar}
\newunicodechar{→}{\ensuremath{\rightarrow}}
% Mismo problema con los simbolos de marca y vinetas: el «✓» de «una palomita ✓»
% salia como cuadrito vacio en el PDF de todas las guias que lo usaban.
\usepackage{pifont}
\newunicodechar{✓}{\ding{51}}
\newunicodechar{✗}{\ding{55}}
\newunicodechar{✔}{\ding{52}}
\newunicodechar{•}{\textbullet}
\newunicodechar{≥}{\ensuremath{\geq}}
\newunicodechar{≤}{\ensuremath{\leq}}

\setmainfont{Helvetica Neue}
\setsansfont{Helvetica Neue}
\setlength{\parindent}{0pt}
\setlength{\parskip}{6pt}
% Sin particion de palabras: en 6.º-7.º una palabra cortada por guion al final
% de linea («Comuni-cacion») frena la lectura mucho mas que una linea corta.
\hyphenpenalty=10000
\exhyphenpenalty=10000
\pretolerance=2000
\tolerance=3000
\setlength{\emergencystretch}{3em}
\linespread{1.10}
\renewcommand{\arraystretch}{1.25}
\setlist{leftmargin=5mm,itemsep=1mm,topsep=1mm}
\newcolumntype{Y}{>{\RaggedRight\arraybackslash}X}

\definecolor{milcMagenta}{HTML}{E6007E}
\definecolor{milcVino}{HTML}{5A0038}
\definecolor{milcTurquesa}{HTML}{0093A5}
\definecolor{milcVerde}{HTML}{58A923}
\definecolor{milcMostaza}{HTML}{E5B400}
\definecolor{milcOcre}{HTML}{B86600}
\definecolor{milcNegro}{HTML}{241816}
\definecolor{milcGris}{HTML}{F7F7F4}
\definecolor{milcLinea}{HTML}{E8E2DD}
\definecolor{milcGradePrimary}{HTML}{0066FF}
\definecolor{milcGradeDark}{HTML}{003D99}
\definecolor{milcGradeSoft}{HTML}{D6E8FF}
\definecolor{milcGradeAccent}{HTML}{CFE7FF}
\definecolor{milcGradeText}{HTML}{FFFFFF}
\color{milcNegro}

\pagestyle{fancy}
\fancyhf{}
\lhead{\small Tecnología e Informática · Grado 11}
\rhead{\small Guía 8 · MILC v3}
\cfoot{\small\thepage}
\renewcommand{\headrulewidth}{0pt}

\newtcolorbox{titlebox}[1]{
  enhanced,colback=#1,colframe=#1,arc=7mm,boxrule=0pt,
  left=7mm,right=7mm,top=3mm,bottom=3mm,
  before skip=8pt,after skip=8pt,
  fontupper=\sffamily\bfseries\Large\color{white}
}

\newtcolorbox{softbox}[3]{
  enhanced,colback=#2,colframe=#1,borderline west={4pt}{0pt}{#1},
  arc=2mm,boxrule=.4pt,left=4mm,right=4mm,top=3mm,bottom=3mm,
  before skip=6pt,after skip=8pt,title={#3},coltitle=white,
  colbacktitle=#1,fonttitle=\sffamily\bfseries,fontupper=\RaggedRight,
  attach boxed title to top left={xshift=3mm,yshift=-2mm},
  boxed title style={arc=2mm, boxrule=0pt}
}

\newtcolorbox{drawbox}[2]{
  enhanced,colback=white,colframe=#1,arc=4mm,boxrule=1pt,
  left=5mm,right=5mm,top=4mm,bottom=4mm,before skip=8pt,after skip=8pt,
  title={#2},coltitle=white,colbacktitle=#1,fonttitle=\sffamily\bfseries,
  fontupper=\RaggedRight,
  attach boxed title to top left={xshift=4mm,yshift=-2mm},
  boxed title style={arc=3mm, boxrule=0pt}
}

\newtcolorbox{infoband}[1]{
  enhanced,colback=#1!8,colframe=#1!50,arc=2mm,boxrule=.5pt,
  left=4mm,right=4mm,top=2mm,bottom=2mm,before skip=4pt,after skip=4pt,
  fontupper=\small\RaggedRight
}

% Puente narrativo entre secciones (contrato editorial Conéctate).
% Da continuidad: "Acabas de X. Ahora vas a Y porque Z."
% Visualmente: banda izquierda en color del grado, texto en itálica pequeña.
\newtcolorbox{puentebox}{
  enhanced,colback=milcGradeSoft,colframe=milcGradeDark,
  borderline west={3pt}{0pt}{milcGradeDark},
  arc=0mm,boxrule=0pt,left=5mm,right=4mm,top=2.5mm,bottom=2.5mm,
  before skip=4pt,after skip=8pt,
  fontupper=\itshape\small\color{milcNegro}\RaggedRight
}
% La flecha va en modo matematico a proposito: Helvetica Neue NO tiene el glifo
% U+2192, asi que \textrightarrow se compilaba a NADA («Missing character: There
% is no → in font Helvetica Neue») y el puente salia sin flecha. En modo
% matematico la flecha viene de la fuente matematica y si se dibuja.
\newcommand{\puente}[1]{%
  \begin{puentebox}%
    \textcolor{milcGradeDark}{$\rightarrow$}\hspace{2mm}#1%
  \end{puentebox}%
}

\newcommand{\linead}[1]{\noindent\color{milcLinea}\rule{#1}{0.5pt}\color{milcNegro}}
\newcommand{\respuesta}[1]{\par\vspace{2mm}\foreach \n in {1,...,#1}{\noindent\color{milcLinea}\rule{\linewidth}{0.5pt}\par\vspace{5mm}}\color{milcNegro}}
\newcommand{\checkbox}{\textcolor{milcGradeDark}{\fbox{\phantom{x}}}\hspace{5pt}}

\newcommand{\phasecard}[4]{
\begin{tcolorbox}[enhanced,colback=#3,colframe=#2,arc=3mm,boxrule=.5pt,height=3.05cm,valign=center,halign=center,left=1.5mm,right=1.5mm,top=2mm,bottom=2mm]
{\sffamily\bfseries\fontsize{22}{24}\selectfont\textcolor{#2}{#1}}\par
{\sffamily\bfseries\small #4}
\end{tcolorbox}
}

% ── Piezas del rediseno 2026 (legibilidad 6.º-11.º) ───────────────────
% Banda de estacion: reemplaza el titlebox macizo. Titulo a la izquierda,
% dato practico (tiempo/modalidad) a la derecha, en una sola linea.
\newcommand{\milcestacion}[2]{%
\begin{tcolorbox}[enhanced,colback=#1,colframe=#1,arc=2mm,boxrule=0pt,
  left=5mm,right=5mm,top=2.6mm,bottom=2.6mm,before skip=11pt,after skip=7pt]
{\sffamily\bfseries\fontsize{15}{18}\selectfont\color{white}#2}
\end{tcolorbox}}

% Tarjeta de paso numerada: sustituye las tablas «Pilar / Aplicacion», que
% eran formato de consulta docente, no de estudiante. El numero va en un
% circulo sobre el borde para que la secuencia se lea de un vistazo.
\newcommand{\milcpaso}[3]{%
\begin{tcolorbox}[enhanced,colback=white,colframe=milcLinea,arc=1.5mm,boxrule=.9pt,
  left=14mm,right=4mm,top=3mm,bottom=3mm,before skip=4pt,after skip=4pt,
  fontupper=\RaggedRight,
  overlay={\node[anchor=north west,circle,fill=#2,text=white,inner sep=0pt,
    minimum size=7.4mm,font=\sffamily\bfseries\small]
    at ([xshift=3.6mm,yshift=-3.2mm]frame.north west) {#1};}]
#3
\end{tcolorbox}}

% Casilla marcable de tamano util con lapiz (la anterior era \fbox{\phantom{x}},
% de alto variable segun el texto que la seguia).
\newcommand{\milcmarca}{\framebox[3.8mm]{\rule{0pt}{3.4mm}}}

% Fila marcable de lista suelta (checks de la estacion 1).
\newcommand{\milccheckrow}[1]{%
\par\vspace{1.8mm}\noindent
\makebox[6.4mm][l]{\raisebox{-.55mm}{\textcolor{milcNegro}{\milcmarca}}}%
\parbox[t]{\dimexpr\linewidth-6.4mm\relax}{\RaggedRight #1}\par}

% Celda marcable dentro de la rubrica (tabla de autoevaluacion). Misma
% sangria francesa, pero pensada para vivir dentro de una columna X.
\newcommand{\milcopcion}[1]{%
\makebox[5.2mm][l]{\raisebox{-.55mm}{\textcolor{milcNegro}{\milcmarca}}}%
\parbox[t]{\dimexpr\linewidth-5.2mm\relax}{\RaggedRight #1}}

% ── Recursos visuales de la guía ──────────────────────────────────────
% Declarados en el bloque `recursos:` del YAML y emitidos por el builder.
% \guiaFigura   → imagen rasterizada o PDF (foto, carta celeste, montaje).
% \guiaDiagrama → fuente TikZ/circuitikz (.tex) incrustada como vector:
%                 es la vía para circuitos y esquemáticos editables.
% #1 = ruta relativa al .tex · #2 = pie (puede ir vacío)
\newcommand{\guiaFigura}[2]{%
\begin{center}
\includegraphics[width=0.88\linewidth,height=0.38\textheight,keepaspectratio]{#1}\\[1.5mm]
{\footnotesize\itshape\textcolor{milcNegro!75}{#2}}
\end{center}
}

\newcommand{\guiaDiagrama}[2]{%
\begin{center}
\input{#1}\\[1.5mm]
{\footnotesize\itshape\textcolor{milcNegro!75}{#2}}
\end{center}
}

\begin{document}
\thispagestyle{empty}

% ── Portada ───────────────────────────────────────────────────────────
% Antes: un rectangulo de color a pagina completa. Se veia bien en pantalla,
% pero en la fotocopiadora del colegio se comia el toner y salia gris sucio.
% Ahora la banda ocupa el tercio superior y el resto va en blanco.
\begin{tikzpicture}[remember picture,overlay]
  \path[fill=milcGradePrimary]
    (current page.north west) rectangle ([yshift=-10.6cm]current page.north east);

  \node[anchor=north west,text=milcGradeText] at ([xshift=2.0cm,yshift=-1.55cm]current page.north west)
    {\sffamily\bfseries\fontsize{12}{14}\selectfont GRADO};
  \node[anchor=north west,text=milcGradeText,opacity=.85] at ([xshift=2.0cm,yshift=-2.35cm]current page.north west)
    {\sffamily\bfseries\fontsize{8.5}{10}\selectfont SERIE GUÍAS · TECNOLOGÍA E INFORMÁTICA};
  \node[anchor=north east,text=milcGradeText,opacity=.85,align=right] at ([xshift=-2.0cm,yshift=-1.55cm]current page.north east)
    {\sffamily\bfseries\fontsize{9}{12}\selectfont I.E. Sor María Juliana\\Cartago, Valle del Cauca};

  \node[anchor=west,text=milcGradeText] (gradonum) at ([xshift=1.85cm,yshift=-7.1cm]current page.north west)
    {\sffamily\bfseries\fontsize{112}{112}\selectfont 11};
  % El circulito de grado (°) solo aplica a las guias de grado («11°»); en
  % Bebras y Semillero el numero grande es un momento/modulo, no un grado.
  % Se posiciona relativo al nodo `gradonum`, no en coordenadas absolutas.
  \draw[milcGradeText,line width=1.6pt]
    ([xshift=2.0mm,yshift=-4.5mm]gradonum.north east) circle (.15cm);

  \node[anchor=west,text=milcGradeText,text width=12.3cm,align=left] at ([xshift=6.5cm,yshift=-6.6cm]current page.north west)
    {
     {\sffamily\bfseries\fontsize{9}{11}\selectfont\MakeUppercase{Período 1 · Presencia y marca digital}}\\[3mm]
     {\sffamily\bfseries\fontsize{21}{25}\selectfont Perfiles profesionales ---\\LinkedIn y portafolio\\coherentes}\\[3.5mm]
     {\sffamily\bfseries\fontsize{15}{18}\selectfont Guía 8-11-TIC}
    };
\end{tikzpicture}

\vspace*{9.4cm}
\vfill

{\sffamily\bfseries\fontsize{8}{10}\selectfont\textcolor{milcGradeDark}{LO QUE TE LLEVAS HOY}}

\begin{tcolorbox}[enhanced,colback=white,colframe=milcNegro,arc=2mm,boxrule=1.6pt,
  left=5mm,right=5mm,top=4mm,bottom=4mm,before skip=3pt,after skip=12pt,
  fontupper=\RaggedRight\fontsize{11}{15.5}\selectfont]
Perfil de LinkedIn completo + portafolio mínimo (3 piezas) enlazado al sitio
\end{tcolorbox}

\begin{tcolorbox}[enhanced,colback=milcGris,colframe=milcGris,arc=2mm,boxrule=0pt,
  left=5mm,right=5mm,top=3.5mm,bottom=3.5mm,after skip=12pt]
{\sffamily\bfseries\fontsize{8}{10}\selectfont\textcolor{milcNegro!55}{ESTA GUÍA ES DE}}\\[3mm]
\begin{tabularx}{\linewidth}{X p{3.2cm} p{3.2cm}}
\linead{\linewidth} & \linead{3.0cm} & \linead{3.0cm}\\[-1mm]
{\fontsize{8.5}{10}\selectfont\textcolor{milcNegro!55}{Nombre completo}} &
{\fontsize{8.5}{10}\selectfont\textcolor{milcNegro!55}{Grupo}} &
{\fontsize{8.5}{10}\selectfont\textcolor{milcNegro!55}{Fecha}}\\
\end{tabularx}
\end{tcolorbox}

\vfill

\noindent\textcolor{milcLinea}{\rule{\linewidth}{1pt}}\\[2mm]
{\sffamily\bfseries\fontsize{9.5}{12}\selectfont Autor: Dr. Álvaro Cárdenas Orozco}\\[1mm]
{\sffamily\fontsize{8.5}{11}\selectfont\textcolor{milcNegro!55}{Metodología MILC · apertura ancestral · pensamiento computacional · triángulo Dussel–estoicismo–Floridi}}

\newpage

\section*{Apertura: Perfiles profesionales --- LinkedIn y portafolio coherentes}

\begin{softbox}{milcGradeDark}{milcGradeSoft}{Saber ancestral}
En el Valle, el Pacífico y los pueblos andinos, las \textbf{``señas''} de identidad abrían o cerraban puertas en segundos. Tu apellido (``los Cárdenas''), tu oficio (``soy hijo del herrero''), tu procedencia (``vengo del Cauca''), tu padrino o tu maestro. Bastaba una conversación de 30 segundos en la plaza para que te ubicaran social y profesionalmente. Las señas eran el \emph{handshake} tradicional: confianza basada en red verificada.
\end{softbox}

\begin{softbox}{milcTurquesa}{milcGris}{Saber contemporáneo}
Hoy ese \emph{handshake} sucede en \textbf{LinkedIn} y en tu \textbf{portafolio público}. Tu foto, tu titular, tu sección ``acerca'', tus experiencias, tus recomendaciones --- todo eso compone tus señas digitales. Cada elemento suma o resta credibilidad. Un perfil incompleto cierra puertas que un perfil sólido abre en segundos.
\end{softbox}

\begin{softbox}{milcMostaza}{milcGris}{Pregunta-puente}
\textit{¿Qué hacían las ``señas'' tradicionales que abrían puertas profesionales en segundos? ¿Qué se conserva, qué cambia, en la versión digital?}
\end{softbox}

\begin{softbox}{milcVerde}{milcGris}{Saber y saber hacer}
Al terminar podrás: (1) \textbf{analizar} 3 perfiles profesionales que admiras y extraer su patrón; (2) \textbf{explicar} qué hace coherentes (o no) los elementos de un perfil; (3) \textbf{crear} tu perfil LinkedIn + portafolio mínimo; (4) \textbf{evaluar} el perfil de un compañero con criterios.
\end{softbox}



\small
\begin{tabularx}{\textwidth}{>{\bfseries}p{3.0cm}Y}
\rowcolor{milcGradePrimary!12}Autor & Dr. Álvaro Cárdenas Orozco\\
DBA & Comunicación profesional en entornos digitales (MEN, Lineamientos T\&I)\\
Referentes & Floridi (infoética) · Dussel (estética de la liberación) · Estoicismo\\
Producto final & Perfil de LinkedIn completo + portafolio mínimo (3 piezas) enlazado al sitio\\
\end{tabularx}
\normalsize

\puente{Antes de empezar, mira el viaje completo. Hoy le das forma a tu \emph{handshake digital} --- el primer contacto profesional con quien no te conoce todavía.}

\milcestacion{milcGradeDark}{Ruta MILC: así vamos a aprender}
\begin{tabularx}{\textwidth}{>{\centering\arraybackslash}X>{\centering\arraybackslash}X>{\centering\arraybackslash}X>{\centering\arraybackslash}X}
\phasecard{1}{milcVerde}{milcGris}{Escucha\\[-1mm]\small Observo, escucho y pregunto.} &
\phasecard{2}{milcTurquesa}{milcGris}{Entender\\[-1mm]\small Sistematizo y ordeno.} &
\phasecard{3}{milcMagenta}{milcGris}{Hacer\\[-1mm]\small Construyo y aplico.} &
\phasecard{4}{milcVino}{milcGris}{Cierre\\[-1mm]\small Evalúo y reflexiono.}
\end{tabularx}


\puente{Antes de armar tu propio LinkedIn, vas a leer otros. Detrás de cada perfil que transmite seriedad hay decisiones coherentes; detrás de los confusos, improvisación.}

\milcestacion{milcVerde}{Estación 1 de 3 · Escucha}

\begin{softbox}{milcVerde}{milcGris}{Escena para escuchar}
\textbf{Actividad 1 · ANALIZA --- 3 perfiles que admiras} (15 min · individual). En LinkedIn busca \textbf{3 personas} cuyo perfil te transmita seriedad y competencia (no necesariamente famosas; profesionales con 3-10 años de carrera). Para cada perfil: ¿cómo es la foto? ¿qué dice su titular? ¿qué tono usa su ``Acerca de''? ¿qué destaca?
\end{softbox}

{\sffamily\bfseries\fontsize{8}{10}\selectfont\textcolor{milcNegro!55}{MARCA CUANDO LO HAYAS HECHO}}
\milccheckrow{Encontré 3 perfiles distintos en mi área de interés.}
\milccheckrow{Observé los 6 elementos clave (foto/titular/sobre/experiencia/destacado/recomendaciones).}
\milccheckrow{Anoté qué tienen en común y qué los diferencia.}

\begin{infoband}{milcVerde}


\textbf{Lo que va al cuaderno (Actividad 1 · ANALIZA).} Título: «Actividad 1 --- 3 perfiles profesionales que admiro». \textbf{Formato:} 3 fichas, una por perfil, con los 6 elementos clave + una columna ``qué transmite''. Al final 1 párrafo de 4 renglones con el patrón común. \textbf{Sabes que terminaste cuando} los 6 elementos están analizados en cada perfil con observaciones específicas (no ``buena foto'').
\end{infoband}

\puente{Ya viste 3 perfiles fuertes. Ahora vamos a entender los \textbf{6 elementos} que tiene todo perfil profesional sólido: foto, titular, sobre mí, experiencia, destacado, recomendaciones.}

\milcestacion{milcTurquesa}{Estación 2 de 3 · Entender}

\begin{softbox}{milcTurquesa}{milcGris}{Lo que vas a entender}
Todo perfil profesional sólido se compone de 6 elementos coherentes entre sí. Vamos a descomponerlos con los pilares del pensamiento computacional para que tu perfil sostenga la promesa que tu marca hace.
\end{softbox}

\milcpaso{1}{milcTurquesa}{\textbf{Descomposición.} Tu perfil se descompone en 6 elementos, cada uno con función específica: (1) \textbf{foto profesional} (de la Guía 7, retrato no selfie); (2) \textbf{titular} (no tu cargo formal, sino tu propuesta de valor de la Guía 1, 220 caracteres máximo); (3) \textbf{sobre mí} (3 párrafos coescritos con IA como en la Guía 5, sin clichés); (4) \textbf{experiencia} (cargos + logros medibles, no descripción de tareas); (5) \textbf{destacado} (3 piezas tangibles de portafolio --- tu sitio de la Guía 4 califica); (6) \textbf{recomendaciones} (1-2 de profesores, mentores o colegas que confirmen tu historia). Sin alguno de los 6, el perfil se siente incompleto.}
\milcpaso{2}{milcTurquesa}{\textbf{Patrones.} Todo perfil fuerte sigue el patrón \textbf{``coherencia transversal''}: lo que dice tu titular se sostiene en tu sobre mí, se demuestra en tu experiencia, se evidencia en tu destacado, y otros lo confirman en tus recomendaciones. Es el principio del ``cinco caras dicen lo mismo''. Si los 6 elementos cuentan historias distintas (titular dice ``creativo'', experiencia dice ``analítico''), la marca se rompe y el reclutador automático pasa de largo en 7 segundos.}
\milcpaso{3}{milcTurquesa}{\textbf{Abstracción.} Lo esencial de tu perfil cabe en una pregunta: \emph{¿después de leer mis 6 elementos, alguien sabe quién soy, qué hago y por qué confiar en mí?}. Si después de leer tu perfil esa persona todavía duda, falta coherencia. El test definitivo: pídele a un compañero que lea tu perfil 30 segundos y luego cierre la pestaña; si puede decirte qué haces en una frase, está bien hecho.}
\milcpaso{4}{milcTurquesa}{\textbf{Algoritmo.} (1) Sube tu foto de retrato profesional (Guía 7); (2) escribe titular con tu propuesta de valor de la Guía 1, 220 caracteres máximo; (3) pega tu ``sobre mí'' coescrito con IA (Guía 5) en el campo ``Acerca de''; (4) llena 1-3 experiencias con logros concretos y medibles (``aumenté X en Y\%'', no ``trabajé en X''); (5) sube 3 piezas a Destacado --- puede ser tu sitio de la Guía 4 más 2 proyectos académicos; (6) escribe a 2 docentes o mentores pidiendo recomendación específica (1 párrafo cada uno, no ``escribe lo que quieras''); (7) publica tu primer post anunciando tu sitio.}

\begin{drawbox}{milcTurquesa}{Anatomía de un perfil profesional coherente}
\textbf{Foto:} retrato profesional, no selfie de fiesta. \\[1mm] \textbf{Titular (220 caracteres):} ``Tu propuesta de valor'' --- no tu cargo. \\[1mm] \textbf{Sobre mí:} tu voz coescrita de la Guía 5, sin clichés. \\[1mm] \textbf{Experiencia:} cada cargo con 3 logros medibles, no descripción de tareas. \\[1mm] \textbf{Destacado:} 3 piezas tangibles (sitio, portafolio, proyecto). \\[1mm] \textbf{Recomendaciones:} 1-2 voces externas que confirmen lo anterior.
\end{drawbox}

\begin{softbox}{milcMostaza}{milcGris}{Errores comunes y su corrección}
(1) Titular tipo cargo (``Estudiante de 11°'', ``Buscando oportunidades''): cierra puertas porque no comunica valor; mejor: tu propuesta de valor de la Guía 1. (2) Sobre mí lleno de clichés (``apasionado por aprender'', ``creativo y proactivo''): se nota a kilómetros que es texto IA sin revisar o copy de plantilla. (3) Experiencia sin logros concretos: solo tareas sin resultados (``responsable de redes sociales'' vs ``aumenté el alcance del grupo de 200 a 800 seguidores en 3 meses''). (4) Sin recomendaciones: tu propio perfil no se valida solo, necesita voces externas que confirmen. (5) Foto inadecuada (selfie de fiesta, filtro agresivo): contradice la formalidad profesional y cierra puertas antes de que el reclutador lea una palabra.
\end{softbox}

\begin{infoband}{milcTurquesa}


\textbf{Lo que va al cuaderno (Actividad 2 · EXPLICA).} Título: «Actividad 2 --- Los 6 elementos de mi perfil». \textbf{Formato:} 6 fichas, una por elemento, cada una con: definición (1 frase tuya), error común a evitar, ejemplo de cómo lo aplicarás a tu propio LinkedIn. \textbf{Sabes que terminaste cuando} cada ficha tiene tu propio ejemplo específico, no plantilla genérica.
\end{infoband}

\puente{Tienes el patrón. Toca aplicarlo a \emph{tu} perfil: foto profesional (de la Guía 7), titular con propuesta de valor (de la Guía 1), texto coherente con tu voz (de la Guía 5). Todo lo que has construido converge aquí.}

\milcestacion{milcMagenta}{Estación 3 de 3 · Hacer}

\begin{softbox}{milcMagenta}{milcGris}{Cómo vas a construir tu producto}
\textbf{Actividad 3 · CREA --- Tu perfil de LinkedIn + portafolio mínimo} (40 min · individual). Hora de armar tu \emph{handshake digital}. Vas a configurar los 6 elementos de tu perfil con todo lo que ya tienes de las guías anteriores: foto, propuesta de valor, voz, sitio. Todo converge aquí.
\end{softbox}

\milcpaso{1}{milcMagenta}{\textbf{Descomposición.} Tu sesión se descompone en 6 pasos --- uno por elemento del perfil. No hagas todo de una vez en bloque caótico: avanza por orden (foto → titular → sobre mí → experiencia → destacado → recomendaciones), revisando cada elemento antes de pasar al siguiente. Esta disciplina pequeña diferencia un perfil ordenado de uno improvisado, y los reclutadores lo notan.}
\milcpaso{2}{milcMagenta}{\textbf{Patrones.} Sigue el patrón \textbf{``crea hoy, refina mañana''}: tu primer perfil no será perfecto, y eso está bien. Lo importante es que esté completo y coherente; ya lo refinarás con feedback en las semanas siguientes. El error opuesto (esperar a tenerlo perfecto para abrirlo) es lo que mantiene al 80\% de los estudiantes sin perfil cuando salen al mercado laboral. Empezar imperfecto vence a esperar perfecto.}
\milcpaso{3}{milcMagenta}{\textbf{Abstracción.} Cada elemento expresa una sola idea de ti. Si tu titular dice ``diseñador'' y tu sobre mí lo contradice diciendo ``ingeniero en formación'', tienes 2 marcas en un solo perfil --- eliges una y ajustas el otro. La coherencia se construye eligiendo qué dejar afuera, no agregando más cosas. Un perfil con 3 áreas distintas es un perfil con 0 marca clara.}
\milcpaso{4}{milcMagenta}{\textbf{Algoritmo.} (1) Entra a LinkedIn o crea cuenta con tu correo institucional; (2) sube foto + titular + sobre mí (15 min); (3) llena 1-3 experiencias con logros concretos y medibles (10 min); (4) sube destacado con enlace a tu sitio de la Guía 4 (5 min); (5) escribe a 2 docentes pidiendo recomendación específica (5 min); (6) publica tu primer post anunciando tu sitio (5 min); (7) revisa el perfil en pestaña incógnita para ver qué muestra a quien no está logueado.}

\begin{drawbox}{milcMagenta}{Checklist antes de declarar el perfil publicado}
\checkbox Foto profesional subida (no selfie). \\[1mm] \checkbox Titular con tu propuesta de valor (no ``estudiante''). \\[1mm] \checkbox Sobre mí coescrito de la Guía 5, sin clichés. \\[1mm] \checkbox 1-3 experiencias con logros concretos, no tareas. \\[1mm] \checkbox 3 piezas en Destacado, incluyendo tu sitio. \\[1mm] \checkbox 2 solicitudes de recomendación enviadas a docentes/mentores.
\end{drawbox}

\begin{softbox}{milcMagenta}{milcGris}{Plantilla de trabajo}
\textbf{Mi titular final.} \textit{220 caracteres máximo, mi propuesta de valor.} \\[1mm] \textbf{Mi sobre mí.} \textit{3 párrafos coescritos de la Guía 5.} \\[1mm] \textbf{Mis 3 logros principales.} \textit{Uno por experiencia, medible.} \\[1mm] \textbf{Mi destacado.} \textit{3 piezas con URL.} \\[1mm] \textbf{Recomendaciones pendientes.} \textit{2 nombres + fechas de envío.}
\end{softbox}

\begin{infoband}{milcMagenta}


\textbf{Lo que va al cuaderno (Actividad 3 · CREA).} Título: «Actividad 3 --- Mi perfil LinkedIn · v1». \textbf{Formato:} 6 secciones con el contenido final de cada elemento + URL pública de tu perfil. \textbf{Extensión:} una página de cuaderno. \textbf{Sabes que terminaste cuando} alguien que abre tu URL entiende en 30 segundos qué haces y siente coherencia entre los 6 elementos.
\end{infoband}

\puente{Lo que sigue es el resumen de \emph{qué} entregas hoy. El criterio principal: que tu LinkedIn y tu portafolio cuenten la misma historia, con la misma voz, sin inflar ni desmerecer.}

\milcestacion{milcMostaza}{Producto de la sesión}
\textbf{Perfil LinkedIn completo + portafolio mínimo enlazado al sitio}.
\begin{softbox}{milcMostaza}{milcGris}{Criterios de calidad}
(1) URL pública del perfil accesible para quien no esté logueado. (2) Los 6 elementos están presentes y completos (no ``estudiante'' como titular). (3) Coherencia transversal: foto + titular + sobre mí + experiencia se sostienen entre sí. (4) Mínimo 3 piezas en Destacado, incluyendo el sitio de la Guía 4. (5) 2 solicitudes de recomendación enviadas (aunque aún no se hayan recibido).
\end{softbox}

\puente{Antes de cerrar, mira tu identidad profesional digital desde las cinco dimensiones humanas. LinkedIn parece técnico pero toca temas profundos de clase, acceso, autenticidad.}

\milcestacion{milcVino}{Evaluación liberadora · cinco dimensiones}

\begin{softbox}{milcVino}{milcGris}{Sentido de la evaluación}
Evaluar no es solo revisar una nota. Es reconocer qué aprendí, qué cambió en mí, qué emociones manejé, cómo me relaciono con la ciudad y la comunidad, y qué le dejaré a quien viene detrás.
\end{softbox}

\textbf{Mi reflexión liberadora en cinco dimensiones.} Completa cada frase iniciada:

\small
\begin{tabularx}{\textwidth}{>{\bfseries\RaggedRight\arraybackslash}p{3.4cm}Y>{\RaggedRight\arraybackslash}p{4.6cm}}
\rowcolor{milcVino}\color{white}Dimensión & \color{white}\bfseries Mi respuesta (frame starter) & \color{white}\bfseries Algo que descubrí\\
1. Personal & Lo que más cambió en mí fue \linead{6.5cm} & \linead{4.2cm}\\
2. Emocional & La emoción más fuerte fue \linead{2.5cm} y la regulé \linead{3cm} & \linead{4.2cm}\\
3. Ciudadana & En democracia esto importa porque \linead{6cm} & \linead{4.2cm}\\
4. Local (Cartago) & En mi barrio sería útil para \linead{6.5cm} & \linead{4.2cm}\\
5. Intergeneracional & A mi abuela/abuelo le diría \linead{6.5cm} & \linead{4.2cm}\\
\end{tabularx}
\normalsize

\begin{infoband}{milcVino}
\textbf{Para la próxima sustentación.} Vuelve a esta tabla en seis meses. Verás que las respuestas cambian — no porque lo aprendido cambie, sino porque tú cambias. La evaluación liberadora es una conversación contigo a través del tiempo.
\end{infoband}


\puente{Y para terminar, tres voces sobre la presencia profesional. Dussel sobre quién está en LinkedIn y quién no, Séneca sobre el tiempo y la reputación, Floridi sobre la identidad como capital reputacional.}

\milcestacion{milcGradeDark}{Triángulo de pensamiento · cierre filosófico}

\begin{softbox}{milcVino}{milcGris}{Dussel · lente del nosotros}
\textit{El sur global piensa desde su propia herida y desde ahí ofrece mundo.}

LinkedIn fue diseñado por y para profesionales del norte global, con un canon de carrera lineal (universidad → primer empleo → ascenso → MBA). Esa narrativa invisibiliza la realidad latinoamericana: trabajos paralelos, emprendimientos informales, oficios familiares heredados. Tu perfil puede reproducir ese canon o puede romperlo nombrando con orgullo lo que en otro lugar ocultarían.

\textbf{¿Mi perfil reproduce un modelo de carrera del norte o reconoce mi historia real del Valle? ¿Qué de mi vida estoy borrando para encajar?}
\end{softbox}
\linead{\linewidth}\\[1mm]
\linead{\linewidth}

\begin{softbox}{milcMostaza}{milcGris}{Estoicismo · Séneca · cuidado interior}
\textit{Mientras esperamos vivir, la vida pasa.}

El error más común con LinkedIn en grado 11 es esperar a tener ``algo que mostrar'' para abrirse el perfil. La reputación profesional se construye con tiempo --- empezar hoy con un perfil honesto y sencillo es ganar 2 años contra quien espera a graduarse para abrirlo.

\textbf{¿Estoy postergando construir mi reputación profesional digital esperando un momento ideal? ¿Qué pierdo cada mes que paso sin perfil?}
\end{softbox}
\linead{\linewidth}\\[1mm]
\linead{\linewidth}

\begin{softbox}{milcTurquesa}{milcGris}{Floridi · lente de la infoesfera}
\textit{La identidad profesional digital es capital reputacional en la infoesfera.}

Tu LinkedIn es activo informacional: aporta señales que otros sistemas (algoritmos de búsqueda, recruiters automáticos, redes de recomendación) usan para ubicarte. Cuanto más coherente y veraz, más valioso es ese capital. Cuanto más inflado, más frágil.

\textbf{¿Mi perfil construye capital reputacional honesto o capital inflado fácil de derrumbar? ¿Qué de mi perfil puedo defender en 5 años?}
\end{softbox}
\linead{\linewidth}\\[1mm]
\linead{\linewidth}

\begin{infoband}{milcNegro}
\textbf{Nota para el docente.} Las tres formulaciones del triángulo son la síntesis del autor de la guía, no citas textuales de los pensadores. Si vas a citarlos en clase, remite a la obra original.
\end{infoband}

\milcestacion{milcVino}{Mi autoevaluación · marca una casilla por fila}

{\small
\setlength{\extrarowheight}{2.2pt}
\begin{tabularx}{\textwidth}{>{\RaggedRight\arraybackslash\bfseries}p{3.0cm}YYY}
\rowcolor{milcVino}\color{white}\bfseries Criterio & \color{white}\bfseries Lo logré & \color{white}\bfseries En proceso & \color{white}\bfseries Necesito apoyo\\[.6mm]
Comprensión del tema & \milcopcion{Explico las ideas con mis palabras.} & \milcopcion{Reconozco las ideas, falta precisión.} & \milcopcion{Necesito repasar lo básico.}\\[.8mm]
\rowcolor{milcGris}Pensamiento computacional & \milcopcion{Usé los pilares al armar mi producto.} & \milcopcion{Usé algunos pilares.} & \milcopcion{Necesito releer la estación 2.}\\[.8mm]
Aplicación y producto & \milcopcion{Mi producto cumple los criterios.} & \milcopcion{Está armado, con ajustes pendientes.} & \milcopcion{Aún está en construcción.}\\[.8mm]
\rowcolor{milcGris}Reflexión 5 dimensiones & \milcopcion{Respondí cada una con honestidad.} & \milcopcion{Respondí algunas con detalle.} & \milcopcion{Mis respuestas son muy generales.}\\[.8mm]
Triángulo de pensamiento & \milcopcion{Conecté las 3 voces con mi trabajo.} & \milcopcion{Conecté al menos una.} & \milcopcion{No las relaciono con mi proceso.}\\[.8mm]
\end{tabularx}}

\vspace{3mm}
\textbf{Hoy entendí que:}\\[1mm]
\linead{\linewidth}\\[1mm]
\linead{\linewidth}

\textbf{Mi compromiso para la próxima sesión es:}\\[1mm]
\linead{\linewidth}\\[1mm]
\linead{\linewidth}

\begin{softbox}{milcMostaza}{milcGris}{Cita de despedida · Marco Aurelio}
\textit{``El obstáculo en el camino se vuelve el camino. Lo que estorba al camino se vuelve camino.''}\\[1mm]
Cada error de hoy, bien observado, se convierte en pieza del camino que pisarás mañana.
\end{softbox}

\small
\begin{tabularx}{\textwidth}{>{\bfseries}p{3.6cm}Y}
\rowcolor{milcGradeDark!12}Nombre completo: & \linead{10cm}\\
Fecha: & \linead{4cm} \quad Curso/grupo: \linead{4cm}\\
Mi firma: & \linead{9cm}\\
\end{tabularx}
\normalsize

\begin{infoband}{milcGradeDark}
\textbf{Para la próxima semana.} Cada día, dedica 10 minutos a interactuar con contenido profesional en LinkedIn (no scrollear; comentar 1 post con criterio y compartir 1 con tu opinión). Esa práctica diaria construye más reputación que un perfil estático impecable. El próximo lunes traes la bitácora de 7 interacciones.
\end{infoband}

\end{document}
